\section{A sharp criterion for the modulus 24}\label{sec:explicit}
For the principal character $\chi_0$ modulo $24$,
\[
 L(s,\chi_0)=\zeta(s)(1-2^{-s})(1-3^{-s}).
\]
The two extra factors have no zeros in $\Re s>0$. In particular, the
zeros in the critical strip and their multiplicities agree. The
smoothed sum gives a quantitative form of this observation. Define
\begin{equation}
 J(x)=\sum_{j\geq1}\frac{x^{-2j}}{2j(2j+1)},\qquad
 C=2+\gamma-\log(4\pi), \label{eq:J}
\end{equation}
and center the eight-spoke sum by
\[
 W_{24}(x)=E_{24}(x)-2\log x-R_{24}(x)-\log(2\pi)+J(x).
\]
\begin{theorem}\label{thm:rh}
The following statements are equivalent.
\begin{enumerate}
\item Every nontrivial zero of $\zeta$ has real part $1/2$.
\item $|W_{24}(x)|\leq C\sqrt{x}$ for every real $x\geq1$.
\item For each $\varepsilon>0$,
$E_{24}(x)=O_\varepsilon(x^{1/2+\varepsilon})$ as $x\to\infty$.
\end{enumerate}
The constant in statement 2 is forced by equality at $x=1$.
\end{theorem}
\begin{proof}
The classical smoothed explicit formula in this normalization is
\begin{equation}
 E(x)=\log(2\pi)-J(x)+
          \sum_\rho\frac{x^\rho}{\rho(1-\rho)}, \label{eq:explicit}
\end{equation}
with one term per nontrivial zero counted with multiplicity. To verify
the constants, invert the Mellin kernel $1/[s(s-1)]$ on $\Re s=2$
against $-(\zeta'/\zeta)(s)x^s$. The inverse kernel gives
$(x/n-1)_+$, including the value zero at $x=n$. Absolute convergence
permits termwise inversion. At $s=1$ the double pole contributes
$x\log x-(1+\gamma)x$. At $s=0$ the residue is
$\zeta'(0)/\zeta(0)=\log(2\pi)$, and the poles at $-2j$ contribute
$-J(x)$. The nontrivial zeros give the sum in \eqref{eq:explicit}.

For completeness, fix an integer $M\geq1$. The good-height estimate
for $\zeta'/\zeta$ gives heights $T_k\to\infty$ for which it is
$O_M(\log^2 T_k)$ uniformly on the horizontal segments from
$-2M-1$ to $2$ at heights $\pm T_k$ \cite[Lemmas 4--5]{kedlaya}.
On $\Re s=-2M-1$, the zeta functional equation and Stirling's formula
bound the logarithmic derivative by $O(\log(2+M+|\Im s|))$,
with an absolute implied constant. The odd real part keeps the
trigonometric term bounded also at height zero. The kernel contributes
$O(T_k^{-2})$ on the horizontal segments. The full integral on the
left line is $O(x^{-2M-1}\log(M+2)/(M+1))$. Taking $k\to\infty$
and then $M\to\infty$ proves \eqref{eq:explicit}.
The classical count $N(T)=O(T\log T)$ gives absolute, locally uniform
convergence of the nontrivial-zero sum for $x\geq1$.

At $x=1$, $E(1)=1+\gamma$ and
$J(1)=\sum_{j\geq1}((2j)^{-1}-(2j+1)^{-1})=1-\log2$.
Hence
\begin{equation}
 \sum_\rho\frac1{\rho(1-\rho)}=C. \label{eq:C}
\end{equation}
This classical evaluation is also recorded by Burnol \cite[p.~3]{burnol}.
Under statement 1, $\rho(1-\rho)=1/4+(\Im\rho)^2>0$ and
$|x^\rho|=\sqrt x$. By Proposition~\ref{prop:local},
$W_{24}$ is precisely the zero sum in \eqref{eq:explicit}. The
triangle inequality and \eqref{eq:C} prove statement 2, with
$W_{24}(1)=C$. Since $0\leq J(x)\leq J(1)$, the correction in
\eqref{eq:correction} is $O(1+\log x)$, so 2 implies 3.

Assume statement 3. Equation \eqref{eq:correction} transfers the same
bound to $E(x)$. Its Mellin transform
$H(s)=\int_1^\infty E(x)x^{-s-1}\dd x$ is then holomorphic on
$\Re s>1/2$, by locally uniform convergence. On $\Re s>1$, direct
integration gives
\[
 H(s)=-\frac{\zeta'(s)}{s(s-1)\zeta(s)}
             -\frac{1}{(s-1)^2}+\frac{1+\gamma}{s-1}.
\]
This identity extends meromorphically to $\Re s>1/2$. A zero
$\rho$ in the right half of the critical strip would give a pole
of the first term with residue $-\operatorname{ord}_\rho(\zeta)/
[\rho(\rho-1)]\ne0$. The other terms are regular there. Holomorphy
rules out this zero, and the functional equation rules out its
reflection in the left half of the strip. This proves statement 1.
\end{proof}
Thus an independent proof of RH by the prime-cross approach would
require an estimate for $E_{24}(x)$ in statement 3, or equivalently
the bound in statement 2. The bounded oscillation in
Theorem~\ref{thm:log} identifies exactly what removing the two Euler
factors contributes to that estimate.
\begin{cor}[A criterion on the integers]\label{cor:limsup}
With $C$ as in \eqref{eq:J},
\[
 \limsup_{n\to\infty}\frac{E_{24}(n)}{\sqrt n}=
 \begin{cases}
 C,&\text{if RH holds},\\
 +\infty,&\text{if RH fails}.
 \end{cases}
\]
In the second case the corresponding lower limit is $-\infty$.
\end{cor}
\begin{proof}
The difference $E_{24}(x)-W_{24}(x)$ is $O(1+\log x)$.
Under RH put $w_\rho=1/(1/4+(\Im\rho)^2)$. Equations
\eqref{eq:explicit} and \eqref{eq:C} give
\[
 e^{-t/2}W_{24}(e^t)=\sum_\rho w_\rho e^{it\Im\rho},
 \qquad \sum_\rho w_\rho=C.
\]
This is Nicolas's zero sum \cite[(1.18)--(1.19), Section~4]{nicolas},
whose real-argument upper limit is $C$. The recurrence argument also
applies to sampling at integers. For any finite set of ordinates and any tolerance, there are
arbitrarily large integer times $t$ at which all their phases are
within that tolerance of $1$. Indeed the phase vectors at positive
integer times have a convergent subsequence in a compact torus.
Differences of sufficiently far separated subsequence terms give the
claimed returns. Truncate the series with tail smaller than a given
$\delta>0$, and apply this argument to its finite part. Replacing
$t$ by $\log\lfloor e^t\rfloor$ changes those finitely many phases
by $o(1)$, while the absolute tail stays below $\delta$. This proves
that the upper limit along integers is $C$.

If RH fails, some zero has real part $\beta>1/2$. Suppose
$E(x)\leq K\sqrt{x}$ eventually. The Mellin transform of the
eventually nonnegative function $K\sqrt{x}-E(x)$ on $[1,\infty)$
is $K/(s-1/2)-H(s)$, initially for $\Re s>1$. Its pole at that zero
forces its abscissa of convergence $\sigma_c$ to satisfy
$\beta\leq\sigma_c\leq1$. Landau's oscillation theorem forces a
singularity at the real point $\sigma_c$ \cite[Theorem 3]{lay}.
But the formula for $H(s)$ in the preceding proof is regular on
the real interval $(1/2,\infty)$. At $1$ this follows by the
Laurent expansion $\zeta(s)=(s-1)^{-1}+\gamma+O(s-1)$, and on
$(1/2,1)$ it follows from the positive alternating eta series and
$1-2^{1-s}\ne0$. This contradiction rules out every such upper
bound. Applying the same argument to $K\sqrt{x}+E(x)$ rules out
every lower bound of the opposite sign.

To pass to integers, on $n<x<n+1$ the slope of $S_1$ is
$\sum_{j\leq n}\Lambda(j)/j=O(\log^2(2n))$, using
$\Lambda(j)\leq\log j$. The other terms of $E$ vary by
$O(\log(2n))$ across the interval. Continuity at each endpoint
gives $E(x)-E(\lfloor x\rfloor)=O(\log^2(2x))$. Together with
\eqref{eq:correction} this transfers both unbounded oscillations to
$E_{24}(n)/\sqrt n$.
\end{proof}
